% ------------------------------------------------------------------
\begin{frame}{今日のゴール / Today's Goal}
  \begin{block}{目標: スティック値を変数に読み取り、演算でモータを個別制御する}
    \begin{itemize}
      \item ESP-NOW 無線通信の仕組みを理解
      \item チャンネル設定で他の受講生との混信を回避
      \item 変数と四則演算でスティック → モータ Duty を計算
      \item オープンループ手動操縦の限界を体感
    \end{itemize}
  \end{block}
\end{frame}

% ------------------------------------------------------------------
\begin{frame}{ESP-NOW 通信フロー / Communication Flow}
  \begin{center}
    \adjustbox{max width=0.95\textwidth, max height=0.82\textheight}{\input{../../_shared/tikz/espnow_dataflow}}
  \end{center}
\end{frame}

% ------------------------------------------------------------------
\begin{frame}{コントローラ各部 / Controller Layout (MODE 3)}
  \begin{center}
    \resizebox{!}{0.80\textheight}{\input{../../_shared/tikz/controller_annotated}}
  \end{center}
\end{frame}

% ------------------------------------------------------------------
\begin{frame}{コントローラ API}
  \begin{table}
    \centering\small
    \begin{tabular}{lll}
      \hline
      \textbf{関数} & \textbf{説明} & \textbf{値域} \\
      \hline
      \api{set\_channel(ch)} & WiFiチャンネル設定 & 1, 6, 11 \\
      \api{rc\_throttle()} & スロットル & 0.0 -- 1.0 \\
      \api{rc\_roll()}     & ロール     & -1.0 -- +1.0 \\
      \api{rc\_pitch()}    & ピッチ     & -1.0 -- +1.0 \\
      \api{rc\_yaw()}      & ヨー       & -1.0 -- +1.0 \\
      \api{is\_armed()}    & ARM状態    & true / false \\
      \api{motor\_set\_duty(id, v)} & モータ個別制御 & id=1--4, v=0.0--1.0 \\
      \hline
    \end{tabular}
  \end{table}
\end{frame}

% ------------------------------------------------------------------
\begin{frame}{ペアリング手順 / Pairing}
  \begin{table}
    \centering\small
    \begin{tabular}{cp{122mm}}
      \hline
      \textbf{Step} & \textbf{操作} \\
      \hline
      1 & コントローラの \textbf{M5ボタンを押しながら電源ON} → LCD に候補一覧画面 ``=== PAIRING ==='' 表示 \\
      2 & StampFly のボタンを \textbf{3秒長押し} → 青LED高速点滅 + ビープ音 \\
      3 & 右スティック上下（または黄ボタン2つ）で自分の機体（MAC下4桁ラベル）に合わせ、画面ボタンで確定 \\
      4 & ペアリング完了 → 機体LEDが緑常灯になり、次回から自動接続 \\
      \hline
    \end{tabular}
  \end{table}

  \begin{alertblock}{注意}
    \small
    \begin{itemize}
      \item 教室で\textbf{複数組が同時に}ペアリングしても、一覧から MAC ラベルが一致する行を選べば取り違えない（近くの他の機体も一覧に出るので要確認）
      \item うまくいかない場合: 両方を再起動して Step 1 からやり直す
    \end{itemize}
  \end{alertblock}
\end{frame}

% ------------------------------------------------------------------
\begin{frame}{チャンネル設定 / Channel Setting}
  教室で複数の StampFly を同時に飛ばすと、同じチャンネル同士で
  \textbf{混信}が起きる。チャンネルを分けて回避する。

  \vspace{0.5em}

  \begin{columns}[c]
    \begin{column}{0.4\textwidth}
      \centering
      \begin{tabular}{cl}
        \hline
        \textbf{Ch} & \textbf{割り当て例} \\
        \hline
        1  & \textcolor{green!70!black}{グループ A（緑）} \\
        6  & \textcolor{yellow!80!black}{グループ B（黄）} \\
        11 & \textcolor{red!80!black}{グループ C（赤）} \\
        \hline
      \end{tabular}
    \end{column}
    \begin{column}{0.55\textwidth}
      \begin{alertblock}{設定方法}
        \api{set\_channel(ch)} を \texttt{setup()} 内で呼ぶ\\
        \textbf{機体とコントローラで同じ番号にすること！}
      \end{alertblock}
    \end{column}
  \end{columns}
\end{frame}

% ------------------------------------------------------------------
\begin{frame}{モータ配置とミキシング / Motor Layout \& Mixing}
  \begin{columns}[T]
    \begin{column}{0.45\textwidth}
      \centering
      \resizebox{0.95\columnwidth}{!}{\input{../../_shared/tikz/motor_layout}}
    \end{column}
    \begin{column}{0.5\textwidth}
      \begin{table}
        \centering\footnotesize
        \begin{tabular}{lcccc}
          \hline
          \textbf{Motor} & \textbf{T} & \textbf{Roll} & \textbf{Pitch} & \textbf{Yaw} \\
          \hline
          M1 FR & + & $-$ & + & + \\
          M2 RR & + & $-$ & $-$ & $-$ \\
          M3 RL & + & + & $-$ & + \\
          M4 FL & + & + & + & $-$ \\
          \hline
        \end{tabular}
      \end{table}

      \vspace{0.5em}

      \begin{exampleblock}{\footnotesize 直感的に理解する}
        \scriptsize
        Pitch+ → 前傾 → 前(M1,M4)強 + 後(M2,M3)弱\\
        Roll+ → 右傾 → 右(M1,M2)弱 + 左(M3,M4)強
      \end{exampleblock}
    \end{column}
  \end{columns}
\end{frame}

% ------------------------------------------------------------------
\begin{frame}{オープンループ制御 / Open-Loop Control}
  \begin{center}
    \resizebox{0.75\textwidth}{!}{\input{../../_shared/tikz/open_loop}}
  \end{center}

  \vspace{0.5em}

  \begin{alertblock}{問題点}
    フィードバックがないため、外乱（風・重心ずれ）に対応できない！\\
    → Lesson 5/8 で PID 制御を導入して解決
  \end{alertblock}
\end{frame}

% ------------------------------------------------------------------
\begin{frame}[fragile]{実習: 手動ミキシング (1/2) / Hands-on: Setup}
  \begin{sflisting}{user\_code.cpp --- setup}
#include "workshop_api.hpp"
static uint32_t tick = 0;
void setup() {
    ws::print("L2: Open-Loop Control");
    ws::arm();              // Enable motor output
    ws::set_channel(1);     // 1, 6, or 11
}
  \end{sflisting}

  \vspace{0.3em}
  \begin{alertblock}{チャンネル番号}
    講師が指定した番号（1 / 6 / 11）に書き換えること！\\
    コントローラも同じチャンネルに設定する
  \end{alertblock}
\end{frame}

% ------------------------------------------------------------------
\begin{frame}[fragile]{実習: 手動ミキシング (2/2) / Hands-on: Loop}
  \begin{sflisting}{user\_code.cpp --- loop}
void loop_400Hz(float dt) {
    tick++;
    float t = ws::rc_throttle();
    float r = ws::rc_roll()  * 0.3f;
    float p = ws::rc_pitch() * 0.3f;
    float y = ws::rc_yaw()   * 0.3f;
    ws::motor_set_duty(1, t - r + p + y); // FR
    ws::motor_set_duty(2, t - r - p - y); // RR
    ws::motor_set_duty(3, t + r - p + y); // RL
    ws::motor_set_duty(4, t + r + p - y); // FL
    if (tick % 80 == 0)
        ws::print("T=%.2f R=%.2f P=%.2f Y=%.2f",
            t, r, p, y);
}
  \end{sflisting}
\end{frame}

% ------------------------------------------------------------------
\begin{frame}{チェックポイント / Checkpoint}
  \begin{alertblock}{確認事項}
    \begin{itemize}
      \item[$\square$] 指定チャンネルで混信なく通信できた
      \item[$\square$] コントローラとペアリングできた
      \item[$\square$] スロットルで全モータが均等に回る
      \item[$\square$] ピッチスティックで前後モータの回転差が出る
    \end{itemize}
  \end{alertblock}

  \vspace{1em}

  \begin{block}{次のレッスン}
    Lesson 3: LED 制御 --- システム状態を LED で可視化
  \end{block}
\end{frame}
